\section*{Capítulo \ref{ch:g7:lungs-and-gas-exchange} --- Los pulmones y el intercambio de gases}

\begin{solution}{exo:g7:lungs-and-gas-exchange:1}
Al tomar aire: los \omterm{def:g3:bones-and-muscles:muscle}{músculos} de las \omterm{def:g3:bones-and-muscles:skeleton}{costillas} levantan y ensanchan la
caja mientras el \omterm{prop:g7:lungs-and-gas-exchange:ventilation}{diafragma} se aplana hacia abajo; el pecho estirado
abre los \omterm{def:g7:how-animals-breathe:four}{pulmones} y el aire entra --- trabajo muscular. Al soltar aire
en reposo: los \omterm{def:g3:bones-and-muscles:muscle}{músculos} se relajan y el pecho estirado recupera su
forma y empuja el aire hacia fuera --- soltar, casi gratis.
\end{solution}

\begin{solution}{exo:g7:lungs-and-gas-exchange:2}
El ancho piso muscular que hay debajo de los \omterm{def:g7:how-animals-breathe:four}{pulmones} y que cierra el
pecho por abajo; su aplanamiento es la mitad profunda de cada toma de
aire.
\end{solution}

\begin{solution}{exo:g7:lungs-and-gas-exchange:3}
Las bolsas de aire de los \omterm{def:g7:how-animals-breathe:four}{pulmones}: cientos de millones de burbujas de
pared fina en las puntas más finas de las vías respiratorias, cada una
envuelta en \omterm{def:g4:heart-and-blood:vessels}{vasos sanguíneos} finos.
\end{solution}

\begin{solution}{exo:g7:lungs-and-gas-exchange:4}
Inspirado: alrededor del $21$ por ciento de \omterm{def:g7:what-is-respiration:respiration}{oxígeno} y casi nada de
\omterm{def:g7:what-is-respiration:respiration}{dióxido de carbono}. Espirado: alrededor del $16$ por ciento de \omterm{def:g7:what-is-respiration:respiration}{oxígeno}
y alrededor del $4$ por ciento de \omterm{def:g7:what-is-respiration:respiration}{dióxido de carbono}. Las diferencias
--- cinco puntos de \omterm{def:g7:what-is-respiration:respiration}{oxígeno} retirados y cuatro de \omterm{def:g7:what-is-respiration:respiration}{dióxido de carbono}
añadidos --- son el intercambio mismo.
\end{solution}

\begin{solution}{exo:g7:lungs-and-gas-exchange:5}
La superficie de paso grande, fina, húmeda y forrada de \omterm{prop:g4:journey-of-food:absorb}{sangre}. La
comparte con las \omterm{def:g1:animals-around-us:covering}{plumas} de las \omterm{def:g7:how-animals-breathe:four}{branquias}
(\cref{prop:g7:how-animals-breathe:spec}) y con las \omterm{prop:g7:digestion-and-nutrients:absorption}{vellosidades} del
intestino (\cref{ex:g7:digestion-and-nutrients:spec}).
\end{solution}

\begin{solution}{exo:g7:lungs-and-gas-exchange:6}
Los alquitranes que recubren e irritan las vías respiratorias (la
escoba sobrecargada y su tos); las paredes alveolares que se deshacen
--- superficie perdida; un gas del humo que ocupa los asientos del
\omterm{def:g7:what-is-respiration:respiration}{oxígeno} en la \omterm{prop:g4:journey-of-food:absorb}{sangre}; y las vías estrechadas que estrangulan el fuelle.
\end{solution}

\begin{solution}{exo:g7:lungs-and-gas-exchange:7}
Tomar aire tiene que estirar el pecho y los \omterm{def:g7:how-animals-breathe:four}{pulmones} contra su muelle
--- ese estiramiento es el trabajo. Soltarlo en reposo es ese mismo
muelle volviendo: el pecho ensanchado y los \omterm{def:g7:how-animals-breathe:four}{pulmones} estirados
recuperan su forma, y ese retroceso empuja el aire.
\end{solution}

\begin{solution}{exo:g7:lungs-and-gas-exchange:8}
El intercambio necesita superficie tanto como aire fresco: el paso solo
ocurre allí donde el aire se encuentra con una pared forrada de \omterm{prop:g4:journey-of-food:absorb}{sangre}.
Un fuelle más rápido a través de bolsas menos numerosas y más flácidas
ventila una superficie que ya no existe --- el factor que falta es la
propia área de paso.
\end{solution}

\begin{solution}{exo:g7:lungs-and-gas-exchange:9}
El fumador sube con las vías estrechadas (toma estrangulada), con una
superficie alveolar reducida (menos carga por \omterm{def:g7:what-is-respiration:respiration}{respiración}) y con
asientos de \omterm{def:g7:what-is-respiration:respiration}{oxígeno} ocupados por el gas del humo (menos transportado
por ronda). Para pagar la misma factura muscular, el \omterm{def:g4:heart-and-blood:heart}{corazón} tiene que
hacer más rondas --- de ahí la llegada jadeando con un \omterm{prop:g4:heart-and-blood:pulse}{pulso} de nivel
de esprint.
\end{solution}

\begin{solution}{exo:g7:lungs-and-gas-exchange:10}
Los dos son control entrenado de la fase de \emph{salida}: le añaden
\omterm{def:g3:bones-and-muscles:muscle}{músculo} medido a lo que normalmente es un muelle pasivo --- el buceador
para ir soltando de manera constante el aire que se expande, y el
trompetista para dosificar una sola salida larga y controlada. (Sus
tomas profundas entrenan también el trabajo de entrada.)
\end{solution}

\begin{solution}{exo:g7:lungs-and-gas-exchange:11}
Una \omterm{def:g7:what-is-respiration:respiration}{ventilación} más profunda y más rápida multiplica la carga en el
mostrador alveolar; un \omterm{def:g4:heart-and-blood:heart}{corazón} más rápido y más fuerte multiplica las
rondas en el mostrador del transporte; y el desvío ensancha los \omterm{def:g4:heart-and-blood:vessels}{vasos}
locales de los \omterm{def:g3:bones-and-muscles:muscle}{músculos} en el mostrador de la entrega. Tres ajustes, un
reparto multiplicado (los libros cuadrados del
\cref{ex:g7:lungs-and-gas-exchange:accounting}).
\end{solution}

\begin{solution}{exo:g7:lungs-and-gas-exchange:12}
La frontera: cientos de millones de \omterm{def:g7:lungs-and-gas-exchange:alveoli}{alvéolos} suman alrededor de media
pista de tenis de superficie de intercambio, plegada dentro de un pecho
--- y por ella cruzan hacia dentro cinco puntos porcentuales del
\omterm{def:g7:what-is-respiration:respiration}{oxígeno} de cada \omterm{def:g7:what-is-respiration:respiration}{respiración} y hacia fuera cuatro de \omterm{def:g7:what-is-respiration:respiration}{dióxido de carbono}.
El fuego lento: los alquitranes y la irritación del humo consumen esa
frontera año tras año --- recubriendo sus accesos, hundiendo sus
bolsas, tomando los asientos de su transporte --- de manera
acumulativa, silenciosa y solo reversible en parte, como una frontera
quemada más deprisa de lo que se puede reconstruir.
\end{solution}

\begin{solution}{pb:g7:lungs-and-gas-exchange:1}
\textbf{1.} Fase de entrada: los \omterm{def:g3:bones-and-muscles:muscle}{músculos} de las \omterm{def:g3:bones-and-muscles:skeleton}{costillas} levantan y
ensanchan la caja y el \omterm{prop:g7:lungs-and-gas-exchange:ventilation}{diafragma} se aplana --- cuesta \omterm{prop:g3:balanced-diet:jobs}{energía}. Fase de
salida (en reposo): los \omterm{def:g3:bones-and-muscles:muscle}{músculos} sueltan y el pecho y los \omterm{def:g7:how-animals-breathe:four}{pulmones}
recuperan su forma --- casi gratis. El ciclo, unas $15$ veces por
minuto.

\textbf{2.} Los \omterm{def:g7:how-animals-breathe:four}{pulmones} se pegan a la pared interior del pecho:
ensancha el pecho y los \omterm{def:g7:how-animals-breathe:four}{pulmones} se estiran con él, y el aire de fuera
entra a llenar el espacio estirado. La bomba funciona cambiando la
forma de la habitación, no apretando el órgano.

\textbf{3.} Unos $15$ ciclos por minuto en reposo; con esfuerzo, los
ciclos se aceleran \emph{y} se hacen más profundos --- más
respiraciones y más aire por \omterm{def:g7:what-is-respiration:respiration}{respiración}.

\textbf{4.} Jadear: ciclos rápidos y superficiales, casi todo
\omterm{prop:g7:lungs-and-gas-exchange:ventilation}{diafragma}. Suspirar: un estiramiento extraprofundo y una soltada larga
--- un reinicio del muelle. Velas: la fase de salida con \omterm{def:g3:bones-and-muscles:muscle}{músculo}
añadido al muelle.

\textbf{5.} Los \omterm{def:g7:lungs-and-gas-exchange:alveoli}{alvéolos}: cientos de millones de burbujas de pared
fina, cada una envuelta en \omterm{def:g4:heart-and-blood:vessels}{vasos} finos, que suman alrededor de media
pista de tenis.

\textbf{6.} Grande: la pista sumada. Fina: paredes de burbuja del
grosor de una \omterm{def:g5:first-look-at-cells:cell}{célula}. Húmeda: su revestimiento. Forrada de \omterm{prop:g4:journey-of-food:absorb}{sangre}: los
\omterm{def:g4:heart-and-blood:vessels}{vasos} que la envuelven. Pliego cumplido en los cuatro puntos.

\textbf{7.} Cinco puntos retenidos de $21$ --- más o menos una cuarta
parte del \omterm{def:g7:what-is-respiration:respiration}{oxígeno} inspirado se retira en cada \omterm{def:g7:what-is-respiration:respiration}{respiración}.

\textbf{8.} De la \omterm{def:g6:exploring-our-environment:conditions}{humedad} de la superficie: el aire que sale ha barrido
unas paredes mojadas (y calientes), así que el escape se lleva agua y
calor --- el espejo empañado es el coste de mantenimiento de la
superficie, pagado en cada \omterm{def:g7:what-is-respiration:respiration}{respiración}.

\textbf{9.} Toma: los alquitranes recubren e irritan las vías. Bomba:
las vías estrechadas estrangulan cada ciclo. Superficie: las paredes
alveolares se deshacen y se funden --- área perdida. Transporte: el gas
del humo ocupa los asientos del \omterm{def:g7:what-is-respiration:respiration}{oxígeno} en la \omterm{prop:g4:journey-of-food:absorb}{sangre}.

\textbf{10.} Las toses y los resfriados inflaman y despejan --- la
maquinaria se repara y vuelve al servicio. Las paredes alveolares rotas
no se reconstruyen: cada bolsa fundida es área de intercambio dada de
baja para siempre --- capital, no gastos corrientes.

\textbf{11.} (1) No admitir humo: las pérdidas de capital son para
siempre. (2) Funcionar con aire limpio: los filtros y las escobas
tienen una capacidad finita. (3) Ejercitar la bomba con carga: el
esfuerzo profundo mantiene los \omterm{def:g3:bones-and-muscles:muscle}{músculos} y ventila la superficie
entera. (4) Tomar el aire por la nariz con frío y con polvo: el
acondicionamiento previo protege el conjunto.

\textbf{12.} La única entre las bombas del cuerpo, el fuelle responde a
la vez a los servicios automáticos y a la voluntad --- porque la misma
salida que ventila el \omterm{def:g7:what-is-respiration:respiration}{dióxido de carbono} es el viento del habla y del
canto, y una máquina prestada al lenguaje tiene que aceptar la mano de
quien habla sobre los mandos.
\end{solution}
